\songtitle{Le cœur léger}

\tag*{2026}\tag{folk-ballad}\tag{french-lyrics}

\begin{notes}
\end{notes}

\begin{style}
Tarantella-leaning instrumental in 6/8 with a fast, galloping pulse; opening hits hard with dungchen, piwang, marimba, fiddle, male and female gaita, double bass, and timbales, then verses ease into a calmer bass-led drive, Pre-chorus swells with layered strings and winds into a bright crescendo, chorus opens into an instrumental break with dancing lead motifs and percussive lift, Outro thins instruments one by one until only timbales remain, crisp and intimate, mix bright, energetic, and folk-cinematic, marimba, ballad
\end{style}

\begin{lyrics}
\songpart{Verse 1}
Dans les vignes du matin\\
Le soleil mord le chemin\\
On ouvre la porte en grand\\
Et le rire sort avant

Le pain chaud sur la table\\
Le vin rouge, incassable\\
Tes yeux brillent, mon ami\\
La fête commence ici

\songpart{Pre-Chorus}
Allez, viens danser\\
La nuit va passer\\
Le cœur léger\\
On va trinquer

\songpart{Chorus}
Ô Languedoc, on vit fort
Ô Languedoc, on rit encore\\
Le verre levé, le cœur d’accord
Ô Languedoc, on vit fort

(Ô Languedoc)
(Ô Languedoc)

\songpart{Verse 2}
Sur la place, les guitares\\
Font tourner les regards\\
Les épaules contre les épaules\\
Et la joie devient parole

La bouteille passe de main\\
Comme un secret, comme un refrain\\
On chante jusqu’au bout du vent\\
Et demain n’existe pas maintenant

\songpart{Pre-Chorus}
Allez, viens danser\\
La nuit va passer\\
Le cœur léger\\
On va trinquer

\songpart{Chorus}
Ô Languedoc, on vit fort
Ô Languedoc, on rit encore\\
Le verre levé, le cœur d’accord
Ô Languedoc, on vit fort

(Ô Languedoc)
(Ô Languedoc)

\songpart{Bridge}
Si le monde est lourd parfois\\
Ici, on le laisse là\\
Une gorgée, un sourire\\
Et tout recommence à vivre

\songpart{Final Chorus}
Ô Languedoc, on vit fort
Ô Languedoc, on rit encore\\
Le verre levé, le cœur d’accord
Ô Languedoc, on vit fort

(Ô Languedoc)
(Ô Languedoc)
\end{lyrics}

